\chapter{Literature Review}

\section{Experimental and First-Principle Methods}

The search for novel solid electrolytes has historically been propelled by two primary approaches: experimental synthesis and characterization, and first-principles computational modeling. Each of these pillars has been instrumental in advancing our understanding of ion transport in solids, yet both face significant limitations that hinder rapid material discovery.

\subsection{Experimental Synthesis and Characterization}

The conventional route to discovering new materials is the empirical, "trial-and-error" approach. This involves the synthesis of a material candidate, followed by its characterization to measure the property of interest. For solid electrolytes, a variety of synthesis techniques are employed, including high-temperature solid-state reactions, sol-gel methods, and co-precipitation, to produce crystalline or amorphous powders which are then sintered into dense pellets for measurement \cite{Goodenough2018}.

The gold standard for measuring ionic conductivity is Electrochemical Impedance Spectroscopy (EIS). In this technique, a small AC voltage is applied across the material over a range of frequencies, and the resulting current is measured. By analyzing the complex impedance response, typically visualized in a Nyquist plot, one can deconvolve the contributions of the bulk material, grain boundaries, and electrode interfaces to the total resistance. From the bulk resistance, the ionic conductivity (\(\sigma\)) can be calculated \cite{Irvine1990}. While indispensable for validating material performance, the entire experimental workflow—from precursor selection and synthesis to pelletization, sintering, and measurement—is inherently slow, laborious, and resource-intensive. Exploring a large compositional space with this method is impractical, creating a major bottleneck in the materials discovery pipeline.

\subsection{First-Principle Computational Methods}

To accelerate the discovery process and gain deeper insights into atomic-scale transport mechanisms, researchers have increasingly turned to first-principles (or \textit{ab initio}) computational methods. These techniques, primarily based on Density Functional Theory (DFT), solve the quantum mechanical equations governing electrons in a material to predict its properties from fundamental physical constants, without any experimental input \cite{Kohn1965}.

DFT calculations can be used to determine the stable crystal structure of a material and to compute the activation energy barrier for ion migration between adjacent sites in the crystal lattice. This energy barrier is a critical parameter that dictates the facility of ion hopping and is directly related to ionic conductivity. For more direct simulations of ion dynamics, methods such as \textit{Ab Initio} Molecular Dynamics (AIMD) are employed. AIMD combines DFT calculations of forces on atoms with classical molecular dynamics, allowing for the explicit simulation of ion trajectories at finite temperatures. From these trajectories, diffusion coefficients and, consequently, ionic conductivities can be calculated \cite{He2017}.

While these first-principles methods offer remarkable accuracy and atomistic insight, their predictive power comes at a steep computational cost. DFT calculations scale cubically with the number of atoms, and AIMD simulations are even more demanding. This restricts their application to relatively small systems (typically a few hundred atoms) and short simulation times (picoseconds to nanoseconds). Consequently, first-principles methods are well-suited for detailed investigations of specific materials or diffusion mechanisms but are too computationally expensive for the high-throughput screening of the thousands or millions of candidate materials required to efficiently explore the vast chemical space. The limitations of both experimental and \textit{ab initio} approaches highlight the need for a new methodology that combines the speed of computation with the ability to screen vast material databases, a gap that machine learning is uniquely positioned to fill.

\section{Machine Learning in Materials Science}

The convergence of large-scale materials data and advanced computational algorithms has given rise to a new paradigm in scientific discovery, often termed the "fourth paradigm" of science: data-driven science \cite{Hey2009}. Within materials science, this has led to the burgeoning field of materials informatics, which leverages machine learning (ML) to accelerate the design, discovery, and deployment of new materials, a core objective of initiatives like the Materials Genome Initiative \cite{MGI2011}.

The fundamental principle of applying ML in this context is to learn a statistical mapping, \(f\), between a material's features (or descriptors), \(\mathbf{x}\), and its target property, \(y\), such that \(y \approx f(\mathbf{x})\). The features, \(\mathbf{x}\), are a set of numerical representations of the material, derived from its composition, structure, or other known attributes. Once trained on a dataset of known materials and their properties, the model \(f\) can make instantaneous predictions for new, un-synthesized candidate materials, bypassing the need for costly experiments or first-principles calculations \cite{Ramprasad2017}. This ability to perform high-throughput virtual screening on vast libraries of potential materials is the primary advantage of the ML approach.

The application of machine learning to the discovery of solid electrolytes is a particularly active and promising area of research. A seminal study by Sendek et al. demonstrated the power of this approach by building an ML model to predict the ionic conductivity of lithium-containing crystalline solids. Their model, trained on data from both experiments and DFT calculations, was used to screen over 12,000 candidate materials, identifying several promising new compositions that were previously overlooked \cite{Sendek2017}. Following this pioneering work, numerous other studies have successfully employed various ML algorithms—from linear models to deep neural networks—to predict ionic conductivity. For example, recent works have focused on developing more sophisticated features and utilizing larger datasets, such as the Database of Doped Solid Electrolytes (DDSE), to improve predictive accuracy and gain deeper insights into the factors governing ion transport \cite{Jalem2021, DDSE2021}.

A critical aspect of building any successful materials informatics model is the process of "feature engineering"—the selection and construction of the descriptors (\(\mathbf{x}\)) that represent the material. These features can be broadly categorized into two types: structure-based and composition-based. Structure-based features encode information about the crystal lattice, such as bond lengths, site coordination, and unit cell volume. While often highly predictive, they require knowledge of the material's crystal structure, which itself can be computationally expensive to obtain. In contrast, composition-based features are derived solely from the stoichiometry and the intrinsic properties of the constituent elements (e.g., atomic mass, electronegativity, ionic radius). As demonstrated in this thesis, these features are particularly powerful for rapid, large-scale screening because they are "structure-agnostic" and can be calculated almost instantaneously for any hypothetical chemical formula, providing a powerful tool to navigate the immense landscape of potential materials.

\section{Research Gap and Problem Statement}

The literature clearly demonstrates a paradigm shift towards data-driven methods for accelerating the discovery of solid electrolytes. While both experimental efforts and first-principles calculations continue to provide essential data and mechanistic insights, they are fundamentally ill-suited for the rapid, large-scale exploration of the vast chemical space required to identify novel, high-conductivity materials. Machine learning has emerged as the most promising strategy to bridge this "discovery gap" by enabling high-throughput virtual screening.

However, a significant challenge within the machine learning approach itself remains. Many of the most accurate predictive models for ionic conductivity rely on features derived from the material's crystal structure \cite{Jalem2021, Sendek2017}. While these structure-based descriptors are powerful, they introduce a critical bottleneck: determining the crystal structure for a hypothetical material is computationally non-trivial. It often requires energy minimization calculations using methods like DFT, which, although faster than full AIMD simulations, are still too slow for screening millions or billions of potential compositions. This reliance on structural inputs limits the speed and scale of virtual screening, constraining the exploration to materials for which a structure is already known or can be readily calculated.

This limitation reveals a distinct research gap: the need for a robust, accurate, and \textbf{structure-agnostic} machine learning model for predicting ionic conductivity. Such a model, based entirely on easily accessible compositional information (i.e., elemental properties), would eliminate the need for any structural pre-computation. This would unlock the ability to perform near-instantaneous screening of virtually any conceivable chemical composition, dramatically expanding the scope of materials discovery.

Therefore, the problem this thesis addresses is the development and validation of a machine learning framework to predict the ionic conductivity of crystalline solids using only elemental-property-based features. The central hypothesis is that a carefully engineered set of compositional descriptors can capture the essential chemical and physical trends that govern ion transport, providing a sufficiently accurate model for a first-pass, high-throughput screening filter. By creating and verifying such a tool, this work aims to provide a computationally inexpensive pathway to down-select promising material families from an enormous compositional space, thereby guiding more focused and efficient experimental and first-principles investigations.
